\newcommand{\drawCube}[3]{
% Define cube coordinates
\coordinate (A#1) at (#2,    0,     0);
\coordinate (B#1) at ({#2+1},0,     0);
\coordinate (C#1) at ({#2+1},1,0);
\coordinate (D#1) at (#2,    1,0);
\coordinate (E#1) at (#2,     0,1);
\coordinate (F#1) at ({#2+1}, 0,1);
\coordinate (G#1) at ({#2+1},1,1);
\coordinate (H#1) at (#2,    1,1);


% Draw cube
\draw[thick] (A#1) -- (B#1) -- (C#1) -- (D#1) -- cycle;
\draw[thick] (E#1) -- (F#1) -- (G#1) -- (H#1) -- cycle;
\draw[thick] (A#1) -- (E#1);
\draw[thick] (B#1) -- (F#1);
\draw[thick] (C#1) -- (G#1);
\draw[thick] (D#1) -- (H#1);
% define lower left corner of doors
\foreach \x in {1,...,#3} {
\pgfmathsetmacro\emptyspace{(1-#3*0.2)/(#3+1)*\x+#2}
\pgfmathsetmacro\doorsspace{0.2*(\x-1)}
\pgfmathsetmacro\llCord{\emptyspace+\doorsspace}
% lower door coordinate
\coordinate (llD\x#1) at (\llCord,0,1);
% Draw doors
\fill[blue!50, rounded corners] (llD\x#1) -- ++(0,0.5,0) -- ++(0.2,0,0)node[pos=0.5, below](in#1\x){\textcolor{blue}{\tiny\texttt{x\x}\normalsize}} -- ++(0,-0.5,0);
}

% coordinate for result door
\coordinate (llDOut#1) at ({#2+1},0,0.2);
\fill[green!50, rounded corners] (llDOut#1) -- ++(0,0.5,0) -- ++(0,0,0.6)node[pos=0.8,rotate=45, below, align=center](out#1){\textcolor{black}{\tiny\texttt{out#1}\normalsize}} -- ++(0,-0.5,0);

% Draw a 2D rectangle around the cube
\node[fit=(E#1) (C#1), inner sep=4pt, red, dotted, draw, line width=1pt,rounded corners](dottedRect#1) {};
% annotate the two-d rect
\node[above=5pt of dottedRect#1.north west, anchor=west, align=left]{\textcolor{red}{\tiny\texttt{def a#1(\textcolor{blue}{\foreach\x in {1,...,#3}{x\x, }...})}\normalsize}};
}

\newtcolorbox{myExBox}[1][]{
  enhanced jigsaw,
  breakable,
  break at=-\baselineskip/0pt,
  parbox=false,
  attach boxed title to top text left={yshift=-3mm,yshifttext=-3mm},
  colback=ForestGreen!10,
  colframe=ForestGreen,
  colbacktitle=ForestGreen,
  #1}

\newtcolorbox{myBox}[1][]{
  enhanced jigsaw,
  parbox=false,
  lines before break=1,
  breakable,
  attach boxed title to top text left={yshift=-3mm,yshifttext=-3mm},
  colback=RoyalBlue!10,
  colframe=RoyalBlue,
  colbacktitle=RoyalBlue,
  #1}
  
\section*{Handout zum Thema ``Anregung von Lernprozessen im Unterricht''}

\begin{myExBox}[title=Beispiel aus der Praxis]
Frustriert trat ich ins Zimmer der Prorektorin. Es ging um die Nachbesprechung meiner Unterrichtssequenz, welche sie soeben hospitiert hatte. Obschon ich mir unglaublich viel Mühe gegeben hatte, kam der Inhalt, welcher ich den SuS vermitteln wollte, leider überhaupt nicht an. Meine Prorektorin kam meine Lektion hospitieren an einem Donnerstagnachmittag, in der zweiten Lektion einer Doppellektion, gerade nachdem die SuS eine Prüfung hatten. Leider waren die SuS, die am Donnerstag einen 10-Lektionen-Tag hatten, fixfertig, und konnten die Inhalte kaum aufnehmen.

Die etwas zu schwer geratene belastete wohl nicht nur die SuS, sondern auch mich, da ich sie ja für Ihre Lern-Aufwände gerecht behandeln wollte. Nun kam jedoch auch noch hinzu, dass die Prorektorin im Raum sass, was dazu führte, dass ich in eine Art ``Notmodus'' schaltete und meine umfangreichen Vorbereitungen demonstrieren und schnell zum Inhalt der Lektion kommen wollte.

Für die Vorbereitung der Inhalte hatte ich mir unglaublich viel Mühe gegeben: Mein Referat zum Thema Funktionen beinhaltete einige Grafiken, für deren Einarbeitung ich bestimmt mehrere Tage gebraucht hatte. Schon nur für die Grafiken hatte ich mehrere Nachmittage investiert, etwa folgende Grafik, die Funktionen mit Abteilungen in einer Fabrik vergleicht:


\begin{figure}[H]
% Use the \drawCube macro to draw various cubes
\centering
\begin{tikzpicture}[scale=2]
\drawCube{1}{0}{2}
% feed input values into a1
\node[fill=blue!50, rounded corners] (input1) at (0.3,0,2){3};
\node[fill=blue!50, rounded corners] (input2) at (0.7,0,2){12};
\draw[thick,->] (input1) -- (in11);
\draw[thick,->] (input2) -- (in12);
\node[fill=green!50, rounded corners](outVal1) at (1.25,-0.6,0.5){36};
\draw[thick,->] (out1) -- (outVal1);
\drawCube{2}{1.75}{3}
\drawCube{3}{3.5}{2}
\draw[thick,->,rounded corners] (out1) -- ++(0,-1.5,0) -- ++(1.05,0,0)node[below]{\lstinline[language=Python,style=TigerJython,basicstyle=\ttfamily\tiny]{return out1}} -- ++(1.55,0,0) -- (in31);
\end{tikzpicture}
\caption*{Selbst erstellte Abbildung, die das Konzept von Funktionen in Python veranschaulichen soll(te)}
\end{figure}
In der Hitze des Gefechts versuchte ich zwar, einige Beispiele und Analogien von Funktionen im echten Leben aufzugeben (Modularität in einer Fabrik, Kochrezepte etc.), allerdings waren es wohl etwas viele Beispiele, und unter Umständen solche, die etwas weit weg vom Leben der SuS schienen. Leider waren nicht nur die SuS wohl nicht mehr wirklich aufnamefähig, sondern auch meine Vorgehensweise, trotz persönlichem Engagement, nicht zielführend. So kam es dazu, dass sich während der Lektion herausstellte, dass ein grosser Teil der SuS leider meinen Erklärungen nicht gefolgt hatte und persönlich nachfragen musste. Dies fiel auch der Prorektorin auf und führte meinerseits zu grosser Frustration, insbesondere angesichts meiner investierten zeitlichen Ressourcen.

\begin{myExBox}[title=Diskussion]
Zentrale theoretische Aspekte, auf die ich eingehen möchte:
\begin{enumerate}
\item Was hätte ich tun können, um den Unterrichtseinstig besser zu gestalten? \cite{kaiser}
\item ...
\item ...
\end{enumerate}
\end{myExBox}
\end{myExBox}
\newpage

\thispagestyle{empty}
\newgeometry{left=.5cm, right=.5cm, top=1cm, bottom=1cm}
\begin{landscape}
\begin{center}
\section*{Literatur: Übersicht \& Zusammenfassung}
\end{center}
\footnotesize
\begin{multicols*}{3}

\begin{myBox}[title=``Unterrichtsplanung und Methoden'' \cite{kaiser}]
Hauptinhalt dieses Texts ist die gestaltung sinnvoller \tib{Unterrichtseinstiege} (UE). Der Text ist folgendermassen gegliedert:
\begin{enumerate}
\item \textbf{Einführung \& Motivation}: Begründung, weshalb dem UE besondere Aufmerksamkeit gewidmet werden soll: er soll 
\begin{enumerate*}
\item neugierig machen
\item das Interesse wecken
\item eine Fragehaltung schaffen
\item zum Kern der neuen Thematik führen 
\item zu einer \tib{Auseinandersetzung} mit dem Sachverhalt führen
\end{enumerate*}.
\item \textbf{Literaturvergleich}: Es wird auf verschiedene Rahmenwerke verwiesen, die sich mit guten UE befasst haben. 
\item \textbf{Systematisierung}: Angesichts der Vielzahl möglicher UE (vgl. didaktische Landkarte) wird im Folgenden eine Systematisierung der UE in drei Zugängen, bzw. Zugriffen auf das Thema vorgenommen:
\begin{enumerate}
\item \textbf{Aggregatszustände}: Dieser Begriff meint, Impulse, Anreize, Anregungen und Interesse via Realbegegnungen, Mediendarstellungen, Modelle, Simulationen, Bericht von ExpertInnen oder Arbeitsbücher zu erwecken.
\item \textbf{Didaktisierungen}: Dieser Begriff meint das didaktische ``Aufladen'' des Unterrichtsgegenstands, um ihn beispielsweise erstaunlicher, bewundernswerter oder bemerkenswerter zu machen. ``Tote'' Gegenstände sollen lebendig gemacht werden: z.B. Gegenstände $\rightarrow$ Erfindungen, Pläne $\rightarrow$ Sorgen, etc., vgl.  ``Roth'sche Rückverwandlung'' \cite{roth}.
\item \textbf{Planungsverfahren}: Grundannahme: Wissen über kommende Inhalte fördert die Lernbereitschaft. Verschiedene Ansätze werden aufgelistet: z.B. Informierender Unterrichtseinstieg \cite{grell}, thematische Landkarte, Themenbörse, Angebotstisch/Themenkiste oder Arbeitsplanarbeit
\end{enumerate}
\end{enumerate}
\end{myBox}\begin{myBox}[title=``Lehr-Lernprozesse in der Schule: Studium'' \cite{maier}]
In diesem Text werden Lernwege im Unterricht (Unterrichtsphasen, Verlaufsformen, sog. Lehr-Lernmodelle) anhand von \tib{Verlaufsmodellen von Unterricht} diskutiert.  \textbf{Verlaufsmodelle} beschreiebn die innere Aufbaulogik einer Unterrichtseinheit, währenddem \textbf{Sichtstrukturen} die Unterrichtsmethoden (Umsetzung des Verlaufsmodells) meinen. Verlaufsmodelle für unterschiedliche Wissensarten werden präsentiert:
\begin{enumerate}
\item \textbf{Direkte Instruktion}: Dieser Ansatz erweist sich nützlich, um einfaches \textbf{deklatratives Wissen} (Faktenwissen) oder einfaches \textbf{prozedurales Wissen} (Regeln, Handlungssequenzen) zu erwerben. Die Phasen dieses Modells sind: \begin{enumerate*}
\item Tägliche Überprüfunge der Hausaufgaben und Übungsaufgeben 
\item neue Lerninhalte präsentieren
\item Übungen durchführen, zu selbstständigen übungen übergehen
\item regelmässige Wiederholung der Fertigkeiten / des Wissens
\end{enumerate*}. Merkmale:  \red{Kleinschrittigkeit}, \red{feedback loops}.
\item \textbf{Concept acquisition}: Um \textbf{konzeptuelles Wissen} zu erwerben, wird das Verlaufsmodell \textit{concept acquisition} präsentiert: 
\begin{enumerate*}
\item konzeptuelles (Vor-)Wissen abrufen
\item Begriffseinführung / Besprecnung
\item Begriffsdefinition (deduktiver Begriffserwerb)
\item Arbeit mit Begriff an weiteren Beispielen
\item Übertragung in neue Kontexte.
\end{enumerate*}. Die \red{Begriffsvernetzung} und mit anderen Begriffen sowie deren \red{Abgrenzung} ist dabei zentral.
\item \textbf{Entdeckendes Lernent}: Dieser Form von Lernen, bei welcher Wissen durch eigene Entdeckungen gewonnen werden soll, wird in der Schule als ``nicht möglich und nicht nötig'' zusammengefasst. Hingegen können ggf. Lernerfolge erzielt werden, indem die SuS mit Fragen der ErfinderInnen des jeweiligen Wissensgegenstands konfrontiert werden und somit die \red{Fragehaltung} übernehmen können. Verlaufsmodelle zum \textit{gelenkten} Entdecken:
\begin{enumerate}
\item \textit{Problem Solving}: \begin{enumerate*}
\item ein Problem wird präsentiert
\item Hypothesen / Vermutungen zur Lösung des Problems werden von SuS entwickelt
\item diese werden überprüft und bewertet
\end{enumerate*}.
\item \textit{Scientific Discovery as Dual Search} (SDDS): 
\begin{enumerate*}
\item Suche im Hypothesenraum
\item Suche im Experimentierraum
\item Analyse von Evidenzen
\end{enumerate*}
\end{enumerate}
\item \textbf{Meisterlehre}: Für den Erwerb von \textbf{komplexem, prozeduralem Wissen} wird inbesondere das Konzept der \textit{coginitive apprenticeship} (\tib{kognitive Meisterlehre)} vergestellt, bei welchem SuS durch Vor- sowie Nachmachen von einem realen Modell (=Meister) lernen. Ein Vorteil davon ist, dass die \red{Bedeutung für das echte Leben} meist klar ist. Verlaufsmodell der Meisterlehre: \begin{enumerate*}
\item \textbf{Modeling} (Vorzeigen)
\item \textbf{Scaffolding} (Unterstützen, Beobachten, Kommentieren)
\item \textbf{Fading} (Schrittweises Abbauen von Unterstützung)
\item \textbf{Coaching} (Hinweise geben, Bewerten der Ausführung)
\end{enumerate*}.
\end{enumerate}
\end{myBox}


\begin{myBox}[title=``Unterrichtsplanung und Methoden'' \cite{gas}]
In diesem Artikel werden zunächste die Vielzähligen Auffassungen von Lernaufgaben präsentiert. \textbf{Lernaufgaben} müssen von \textbf{Prüfungs-} oder \textbf{Leistungsaufgaben} abgegrenzt werden. Am passendensten scheint die Definition als ``jedwede Aufgabenstellungen, die das Lernen der SuS unterstützen''. Ein Kategorisierungssystem für kompetenzfördernde Aufgabenstellugen in mehreren Dimensionen wird vorgestellt \cite{chem}.\vspace{-.3cm}
\begin{table}[H]
\footnotesize
\centering
\begin{tabular}{l|l}
Merkmalsbereich & Merkmal \\\midrule\midrule
Authentizität & \begin{tabular}[c]{@{}l@{}}- Kompetenzabbild\\ - Lebensnähe\end{tabular} \\\midrule
Kognition & \begin{tabular}[c]{@{}l@{}}- Arbeit an (Prä-)Konzepten\\ - Art des Wissens\\ - Kognitiver Prozess\end{tabular} \\\midrule
Komplexität & \begin{tabular}[c]{@{}l@{}}- Strukturierung\\ - Repräsentationsform\end{tabular} \\\midrule
Differenzierung & \begin{tabular}[c]{@{}l@{}}- Offenheit\\ - Lernunterstützung\\ - Vielfalt der Lernwege\end{tabular}
\end{tabular}
\end{table}\vspace{-.5cm}
Die Lernrelevanten Merkmale können jeweils in unterschiedliche Ausprägungen weiter aufgeteilt werden.
\end{myBox}
\begin{myBox}[title=``Gestaltung von Unterricht'' \cite{tulo}]
In einem ersten Teil werden Kriterien an lernprozessförderliche Aufgaben eingeführt und aufgelistet: 
\begin{enumerate*}
\item \textbf{Verständlichkeit} (auf den Erfahrungs- und Vorstellungshorizont der SuS bezogen)
\item \textbf{Situierung} (hinreichend komplex, um anwendungsfähiges Wissen zu generieren)
\item \textbf{Bedeutsamkeit} (Bedürfnisse und Interessen der SuS ansprechen)
\item \textbf{Neuigkeitswert}
\item \textbf{Lösbarkeit}
\item \textbf{Exemplarischer Charakter}
\end{enumerate*}.
Vier komplexe Aufgabenarten zur Anregung von Lernprozessen werden danach vorgestellt und anhand unterschiedlicher Beispiele illustriert:
\begin{enumerate}
\item \textbf{Komplexe Probleme}: Wichtig bei diesem Aufgabentyp ist, dass es eine Lösung im Sinne einer Antwort ``richtig / falsch'' gibt, um die Aufgabe in eine Ausgangslage, einen Transformationsprozess und einen Ziel- bzw. Endzustand (Problemlösung) aufzugliedern und um die Aufgabe von den anderen drei Aufgabentypen abzugrenzen. Dieser Aufgabetyp scheint insbesondere Geeignet für die Erarbeitung von Regeln, Abläufen, Gesetze, Verfahren, Begriffen und Konzepten.
\item \textbf{Komplexe Entscheidungsfälle}: Hier müssen Lernende aus mehreren Entscheidungsmöglichkeiten eine auswählen, ohne dass eine Option aus empirischer oder logischer Sicht klar richtig ist. Es kann unterschieden werden in Fälle, bei denen es einen Ziel- einen Wegkonflikt oder beides gibt.
\item \textbf{Komplexe Gestaltungsaufgaben}: Aufgaben, bei denen etwas gedanklich entworfen und gestaltet werden muss.
\item \textbf{Komplexe Beurteilungen}: Gewissermassen den anderen Weg herum, eine bereits gelöste Aufgabe der anderen drei Typen wird analysiert und bewertet.
\end{enumerate}
Gemeinsam ist all diesen Aufgaben dass Sie nebste dem Ziel, den Lerngegenstand den Lernenden zu erschliessen, Beiträge zur allgemeinen Intellektuellen und sozial-moralischen Entwicklung leisten. Auch sollten entsprechende Aufgaben zu \red{Beginn einer Unterrichtssequenz stehen}, um die Motivation zu steigern und die Sinnhaftigkeit bewusst zu machen. Die eingangs genannten Kriterien können dazu dienen, existierende Aufgaben zu vergleichen auf Verbesserungsmöglichkeiten hin zu analysieren (z.B. für eigenen Unterricht oder Hospitation).
\end{myBox} 
\end{multicols*}
\end{landscape}
\newpage
\restoregeometry
